\scanpage{169}
\mspage{43}

$\unsure{[}$Ainsi, le seul problème, une fois trouvé $C_0$
(petite) satisfaisant $C_0\subset\mathrm{Cof}_W,\ C_{0*}\subset W$,
c'est de trouver une \uuline{petite} $TC_0\subset TC$
qui ``engendre'' $TC$``$=C\cap W$'', i.e. qui satisfasse
aux conditions équivalentes a, a', b, c de la
proposition précédente. J'ai passé des jours
à essayer de prouver l'existence de $TC_0$,
en supposant $\unsure{\mathcal M|M}$ muni d'une filtration
\unsure{cardinal|cardinale}, \inserted{et} $W$ une partie accessible de $\mathcal M$,
sans y être arrivé. Bien sûr, \deleted{\unsure{tout}} ce qui
nous importe vraiment, c'est que $(TC,F)$ satisfasse
à la condition de factorisation. Mais je
ne vois aucun moyen de le prouver, si ce
n'est en ``engendrant'' \inserted{à g.} $TC$ \deleted{\unsure{à g.}} par une \uline{petite}
\unsure{sous-ensemble} $TC_0$.

5. \uline{Résultats techniques préliminaires, sur les parties
accessibles de $\mathrm{Ob}\,\mathcal M$, stables par $\underrightarrow{\mathrm{lim}}$ filtrantes}

Je crois pourtant que je peux arriver à
quelque chose d'utilisable, en travaillant systématiquement \unsure{avec|sur} des $C$ ($\subset\mathrm{Cof}_W$, Q-closes à
gauche, Q-engendrées \inserted{à g.} par une \uline{petite} partie
$C_0$ \deleted{\unsure{)}} de $C$) qui sont stables par \unsure{ttes|toutes} limites
inductives filtrantes.

Je suppose à nouveau \deleted{\unsure{\mathcal M}} sur $\mathcal M$ les
conditions suivantes :

% archive sheet number: 168
